\section*{Methods Source Reference}

This note records the two source documents read on 2026-06-30 for writing the manuscript Methods section. The files below remain the source of truth; this note is a compact index of the details most likely to be needed in the paper.

\subsection*{Synthetic Target-Rescaling Test}

Source:
\begin{verbatim}
Target_Design/monolithic_mechanism/1d_gaussian_delta_normalized_targets/
docs/scaled_target_rescaling_methods.md
\end{verbatim}

Purpose: test whether reducing target magnitude or target variance alone makes a learning problem intrinsically easier. The experiment constructs scaled versions of the same one-dimensional target:
\[
  y(x), \quad 0.75y(x), \quad 0.5y(x), \quad 0.25y(x).
\]
Scaling changes only the amplitude; the waveform, oscillations, local roughness, peaks, and valleys remain unchanged.

Synthetic target:
\[
  x \in [-4,4], \qquad y(x) = g(x) + \Delta(x).
\]
The baseline component is a Gaussian with \(A_g=1.0\), \(\mu=0\), and \(\sigma=1.15\). The residual component is an envelope-modulated combination of sinusoidal and localized step-like terms, mean-centered after construction. The full formula is in the source document.

Data splits:
\begin{itemize}
  \item Train: 50,000 points, seed 1729.
  \item Validation: 50,000 points, seed 1730.
  \item Test: 50,000 points, seed 1731.
  \item All splits sample \(x \sim \mathrm{Uniform}(-4,4)\); this is an in-domain scaling control, not an extrapolation test.
\end{itemize}

Model and training:
\begin{itemize}
  \item Same scalar-input, scalar-output MLP trained from scratch for each scaled target.
  \item Architecture: \(x \rightarrow 64 \rightarrow 64 \rightarrow 1\), SiLU hidden activations, linear output.
  \item Trainable parameters: 4,353.
  \item Target normalization uses each target's train-set mean and standard deviation.
  \item Predictions are transformed back to raw target units before MAE/RMSE.
  \item Optimizer Adam, learning rate \(10^{-3}\), 500 epochs, batch size 1024, seed 0, CPU recorded run, MSE loss on normalized target.
\end{itemize}

Metrics:
\[
  \mathrm{MAE} = \frac{1}{N}\sum_i |\hat{y}_{c,i} - y_{c,i}|,
  \qquad
  n\text{-MAE} = \frac{\mathrm{MAE}}{\mathrm{IQR}(y_c)}.
\]
The IQR is computed separately on each scaled validation target. Validation MAE decreases approximately linearly with scale, while MAE/IQR remains essentially constant.

Figure and reproducibility:
\begin{verbatim}
Target_Design/monolithic_mechanism/1d_gaussian_delta_normalized_targets/
figures/publication_scaled_target_rescaling/scaled_target_rescaling_bipanel.png

Target_Design/monolithic_mechanism/1d_gaussian_delta_normalized_targets/
runs/main_1seed_500epoch_normalized_targets_cpu/
normalized_target_final_metrics.csv
\end{verbatim}

\subsection*{Targets, Baselines, and GNN Methods}

Source:
\begin{verbatim}
New_Analysis/expanded_target_panel/TARGETS_MODELS_METHODS.md
\end{verbatim}

Scope: tmQM target generation, reference corrections, BoE corrections, descriptor baseline targets, MLP baseline targets, GemNet/SchNet target LMDB preparation, training protocols, split-aware preprocessing, and validation rules.

Dataset and splits:
\begin{verbatim}
trillium_project/split_ids/tmqm_train_ids.csv
trillium_project/split_ids/tmqm_val_ids.csv
trillium_project/split_ids/tmqm_test_ids.csv
\end{verbatim}
Fully aligned DFT/PRE/PRE+A target split sizes:
\[
  n_{\rm train}=56830,\qquad n_{\rm val}=7104,\qquad n_{\rm test}=7104.
\]

Reference correction:
\[
  E_{\rm ref}^{m}(i)
  =
  E_{\rm raw}^{m}(i)
  -
  \sum_Z n_Z(i)\epsilon_Z^{m},
\]
where \(m\) is DFT or xTB and atomic references are method-specific.

BoE correction:
\[
  E_{\rm ref,BoE}^{m}(i)
  =
  E_{\rm ref}^{m}(i)
  -
  \mathrm{BoE\_fit}_m(i).
\]
The xTB residual is defined after DFT and xTB each receive their own reference and BoE corrections:
\[
  \Delta^{\rm xTB}_{\rm ref,BoE}
  =
  E_{\rm ref,BoE}^{\rm DFT}
  -
  E_{\rm ref,BoE}^{\rm xTB}.
\]

Canonical target panel:
\begin{itemize}
  \item \(E_{\rm ref,BoE}^{\rm DFT}\): `DFT_ref,BoE`.
  \item \(E_{\rm ref,BoE}^{\rm xTB}\): `xTB_ref,BoE`.
  \item \(\Delta^{\rm xTB}_{\rm ref,BoE}\).
  \item \(\Delta^{\rm BoB}_{\rm ref,BoE}\).
  \item \(\Delta^{\rm PRE}_{\rm ref,BoE}\).
  \item `BoB_ref,BoE`, `PRE_ref,BoE`, `PRE+A_ref,BoE`.
  \item \(\Delta^{\rm PRE+A}_{\rm ref,BoE}\).
\end{itemize}
Registry:
\begin{verbatim}
New_Analysis/expanded_target_panel/target_registry.csv
\end{verbatim}

Descriptor baselines:
\begin{itemize}
  \item BoB: Bag-of-Bonds descriptor, 4,950 dimensions.
  \item PRE: pair-radial expansion descriptor, pair-only typed radial features, 34,650 dimensions.
  \item PRE+A: PRE plus compact central-element angular histogram, 45,738 total dimensions.
\end{itemize}
PRE+A angular defaults: cutoff 6.0 Angstrom, 6 radial shells, 21 radial shell pairs, 12 angle bins uniform in \(\cos\theta\), central element vocabulary size 44, angular dimensions \(44 \times 21 \times 12 = 11088\).

Linear ridge baselines:
\begin{itemize}
  \item Main script: `New_Analysis/expanded_target_panel/run_expanded_ridge_targets.py`.
  \item Split-aware preprocessing only: drop zero-variance channels on train, scale by train-set standard deviation, no centering for sparse matrices, train-only 5-fold CV for alpha, validation/test never used in filtering/scaling/CV/coefficient fitting.
  \item Default CV folds 5, shuffle true, random state 42, zero variance threshold \(10^{-12}\), sklearn Ridge, alpha selected by CV MAE.
\end{itemize}

MLP baseline targets:
\begin{itemize}
  \item Main script: `New_Analysis/mlp_baselines/train_mlp_baselines.py`.
  \item MLP baselines fit BoB, PRE, or PRE+A to `DFT_ref,BoE`; residual targets subtract these fitted baselines.
  \item Architecture: hidden dimensions 512 and 128, SiLU, dropout 0.10, scalar linear output.
  \item Optimizer AdamW, learning rate \(10^{-3}\), weight decay \(10^{-4}\), SmoothL1Loss beta 0.10, batch size 1024, max epochs 800, patience 80, seed 101.
\end{itemize}

GemNet:
\begin{itemize}
  \item Main high-capacity local geometric GNN.
  \item Representative model: `gemnet_t`, 3 blocks, atom/edge embedding size 512, triplet embedding size 64, cutoff 6.0 Angstrom, max neighbors 50, Gaussian radial basis, spherical-harmonic angular basis, SiLU activation, extensive sum, no forces, no PBC.
  \item Local by construction: graph uses neighbors within cutoff; explicit features are pair distances and triplet geometry; higher-order/longer-range effects are approximated through stacked message passing.
  \item Training uses FairChem LMDBs with train-set target normalization.
  \item Validation MAE is taken from the best validation checkpoint in training logs.
\end{itemize}

SchNet:
\begin{itemize}
  \item Controlled lower-capacity local geometric GNN.
  \item Config: hidden channels 128, filters 128, interactions 6, Gaussians 50, cutoff 6.0 Angstrom, sum readout, no forces, no PBC.
  \item Same train-set target normalization principle.
\end{itemize}

Target LMDB generation:
\begin{itemize}
  \item Reads canonical target CSVs by `mol_id`.
  \item Replaces each molecule's energy target in source tmQM LMDB split.
  \item Writes target-specific train/val/test LMDBs.
  \item Computes train-set target mean/std and writes target-specific FairChem configs.
\end{itemize}

Validation and reproducibility rules:
\begin{itemize}
  \item All target CSVs keyed by `mol_id`.
  \item Fixed GemNet split IDs define train/val/test rows.
  \item Baseline preprocessing fit on train only.
  \item Validation/test rows never enter filtering, scaling, alpha CV, or coefficient fitting.
  \item GNN target LMDBs match source LMDB molecules to target CSV rows by `mol_id`.
  \item Target-specific LMDBs checked by preflight scripts.
  \item Reported normalized MAE uses validation target IQR.
\end{itemize}

Manuscript framing from source: the baseline descriptors define chemically motivated target transformations, but the primary predictive models are GemNet and SchNet. The core question is how well GemNet and SchNet learn each transformed target on the same tmQM split, not a competition among all possible baseline regressors.
